% =============================================================================
% 06-continuidade.tex - PARTE 6, o depois. Frames 36 a 38.
% =============================================================================

% ----------------------------------------------------------------- frame 36 --
\begin{frame}{\ftit{\faFile}{Notifica\c{c}\~ao de viol\^encia autoprovocada}}
\CRar{0.8}
\begin{caixanorma}{A obriga\c{c}\~ao}
Casos suspeitos ou confirmados de viol\^encia autoprovocada s\~ao de
notifica\c{c}\~ao compuls\'oria pelos estabelecimentos de sa\'ude, p\'ublicos
e privados, \`as autoridades sanit\'arias. A notifica\c{c}\~ao \'e sigilosa.
\end{caixanorma}

\CRvg
\begin{columns}[onlytextwidth, t]
\begin{column}{0.47\textwidth}
\begin{caixanorma}{O que \'e evento notific\'avel}
\begin{itemc}
\item tentativa de suic\'idio
\item automutila\c{c}\~ao, com ou sem idea\c{c}\~ao
\item suic\'idio consumado
\end{itemc}
\vspace{0.06cm}
\norma{Idea\c{c}\~ao isolada} n\~ao \'e, por si, evento notific\'avel.
\end{caixanorma}
\end{column}
\calha
\begin{column}{0.47\textwidth}
\begin{caixanorma}{Como se faz}
\begin{itemc}
\item Ficha de Notifica\c{c}\~ao de Viol\^encia Interpessoal\slash Autoprovocada
\item lan\c{c}amento no SINAN
\item tentativa de suic\'idio \'e de notifica\c{c}\~ao \norma{imediata, em
  at\'e 24 horas}
\end{itemc}
\end{caixanorma}
\end{column}
\end{columns}

\mensagem{Notificar n\~ao \'e denunciar.}
\fonte{Lei 13.819/2019, art. 6; Portaria GM\slash MS 1.061/2020, item 48.
\selonorma}
\end{frame}

% ----------------------------------------------------------------- frame 37 --
\begin{frame}{\ftit{\faNotesMedical}{Registro cl\'inico}}
\CRar{5}
\begin{columns}[onlytextwidth, t]
\begin{column}{0.47\textwidth}
\begin{caixaavaliar}{Fun\c{c}\~ao cl\'inica}
Um bom registro permite que outra pessoa, no pr\'oximo plant\~ao, reconstrua a
decis\~ao cl\'inica e saiba o que j\'a foi tentado, o que foi combinado e com
quem.
\end{caixaavaliar}
\end{column}
\calha
\begin{column}{0.47\textwidth}
\begin{caixanorma}{Fun\c{c}\~ao m\'edico-legal}
O mesmo registro \'e o que sustenta a conduta em eventual questionamento. N\~ao
se escreve para o processo; escreve-se bem, e isso serve para os dois.
\end{caixanorma}
\end{column}
\end{columns}

\CRvg
\begin{center}
\begin{tikzpicture}[node distance=0.18cm]
  \node[crnode={campoazul}{9.0cm}, font=\scriptsize, inner ysep=2.8pt] (a)
    {\textbf{1.} o que foi identificado, com as falas que importam};
  \node[crnode={campoazul}{9.0cm}, font=\scriptsize, inner ysep=2.8pt, below=of a] (b)
    {\textbf{2.} a formula\c{c}\~ao constru\'ida, e n\~ao a categoria};
  \node[crnode={campobordo}{9.0cm}, font=\scriptsize, inner ysep=2.8pt, below=of b] (c)
    {\textbf{3.} por que a disposi\c{c}\~ao escolhida, e n\~ao a outra};
  \node[crnode={campoverde}{9.0cm}, font=\scriptsize, inner ysep=2.8pt, below=of c] (d)
    {\textbf{4.} as medidas concretas organizadas, e quem participou};
  \node[crnode={campoverde}{9.0cm}, font=\scriptsize, inner ysep=2.8pt, below=of d] (e)
    {\textbf{5.} quando, onde e com quem ser\'a o seguimento};
\end{tikzpicture}
\end{center}
\end{frame}

% ----------------------------------------------------------------- frame 38 --
\begin{frame}{\ftit{\faLink}{Continuidade}}
\CRar{3.7}
\begin{center}
\begin{tikzpicture}
  \node[musialGreen, text width=9.2cm, font=\small, inner ysep=7pt] {%
    \sffamily\bfseries A alta n\~ao encerra o risco. Ela transfere o cuidado para
    outro contexto.};
\end{tikzpicture}
\end{center}

\CRvg
\begin{columns}[onlytextwidth, t]
\begin{column}{0.47\textwidth}
\begin{caixamudanca}{O per\'iodo mais delicado}
Os dias que seguem uma alta ou uma passagem pela urg\^encia est\~ao entre as
janelas de vigil\^ancia. \'E quando o servi\c{c}o presume que o caso se
resolveu e o paciente presume que incomodou.
\end{caixamudanca}
\end{column}
\calha
\begin{column}{0.47\textwidth}
\begin{caixaproteger}{O que sustenta a transfer\^encia}
\begin{itemc}
\item destino com nome, endere\c{c}o e data
\item algu\'em avisado do outro lado
\item contato ativo nas primeiras 48 a 72 horas
\item a fam\'ilia sabendo o que fazer, e n\~ao s\'o o que evitar
\end{itemc}
\end{caixaproteger}
\end{column}
\end{columns}
\end{frame}
